% TOP_PROBLEM: {"id": 8700035, "title": "Conjecture 3.17", "category_id": 3, "category": {"id": 3, "name": "graph_theory", "display_name": "Graph Theory", "description": "Problems involving graphs, networks, and their properties.", "slug": "graph-theory", "order_index": 3, "created_at": "2026-07-31T15:26:25.671Z"}, "status": "open", "problem_number": "AMR-086-0035", "set_id": 13}
% FIRST_POSTED: 2026-10-01
% ENTRY_KIND: statement_only
% UnsolvedMath ID 8700035; source catalogue code AMR-086-0035
% !TeX program = lualatex
% Definition and sources only; this file is not an LLM solution attempt.
\documentclass[11pt,a4paper]{article}
\usepackage[margin=25mm]{geometry}
\usepackage{fontspec}
\setmainfont{lmroman10-regular.otf}[BoldFont=lmroman10-bold.otf,
 ItalicFont=lmroman10-italic.otf,BoldItalicFont=lmroman10-bolditalic.otf]
\usepackage{amsmath,amssymb,mathtools}
\usepackage{unicode-math}
\setmathfont{latinmodern-math.otf}
\AtBeginDocument{\let\setminus\smallsetminus}
\usepackage{xurl,hyperref}
\hypersetup{colorlinks=true,urlcolor=blue}
\setcounter{secnumdepth}{0}
\setlength{\parindent}{0pt}
\setlength{\parskip}{5pt}
\setlength{\emergencystretch}{3em}
\begin{document}
\section*{Conjecture 3.17}\label{8700035}
{\small UnsolvedMath ID 8700035 · AMR-086-0035 · Graph Theory · AMR Open Problem Lists}
\subsection{Definitions and mathematical statement}
For real $s>1$, the Riemann zeta value is defined by the absolutely convergent series
\[
\zeta(s)=\sum_{m=1}^{\infty}m^{-s}.
\]
Consider the infinite list consisting of the usual circle constant $\pi$ and the odd integer zeta values $\zeta(3),\zeta(5),\zeta(7),\ldots$.

Are these numbers algebraically independent over $\mathbb Q$? For an infinite family, this means that every finite subfamily is algebraically independent. Equivalently, for every integer $r\ge1$ and every nonzero polynomial $F\in\mathbb Q[X_0,X_1,\ldots,X_r]$, does one have
\[
F\bigl(\pi,\zeta(3),\zeta(5),\ldots,\zeta(2r+1)\bigr)\ne0?
\]
Polynomials may omit any of their variables, so this finite-prefix formulation includes all possible finite subfamilies and all individual transcendence assertions.

The coefficients must be rational, but clearing denominators gives the equivalent formulation with integer coefficients. The conclusion excludes nonlinear relations of arbitrary degree, not just rational linear relations or relations within some chosen weight.

Only odd positive integer arguments greater than one appear among the zeta values. Even zeta values have known expressions in rational multiples of powers of $\pi$, so including them would create polynomial relations. The question also omits the divergent series at $s=1$ and does not involve analytic continuation. Its precise target is simultaneous algebraic freedom of the displayed real constants, uniformly over every finite number of them.
\subsection{Short English statement}
Are pi and all odd zeta values at integers greater than one jointly algebraically independent over the rationals?
\subsection{Sources}
\begin{itemize}
\item[S1] ULAM.AI and the UnsolvedMath Contributors, supplied problems.json, record 8700035 (AMR-086-0035), AMR Open Problem Lists. Catalogue identity and source transcription; CC BY 4.0.
\url{https://huggingface.co/datasets/ulamai/UnsolvedMath}.
\item[S2] Waldschmidt - Open Diophantine Problems (2004), Conjecture 3.17. Original source indexed by AMR-086-0035.
\url{https://arxiv.org/abs/math/0312440}.
\end{itemize}
\end{document}
